% 【深论模块】所属：第6章（由 chapters/revised/06-formal.tex \input 引入）
\minitopic{一个后验数不等于一个稳健结论}开篇检测得到的$15.4\%$，只是在一组明确假设下的结果：基础率是$1\%$，阳性的似然比为18，竞争假设只有“患病”和“未患病”，检测参数也被当作已知。改变其中任何一项，后验都可能移动。最简单的敏感性分析不是再给更多小数，而是把几个合理设定并列：

\begin{center}
\begin{tabular}{ccc}
\toprule
患病基础率 & 阳性似然比 & 阳性后的患病概率\\
\midrule
$1\%$ & 18 & $15.4\%$\\
$2\%$ & 18 & $26.9\%$\\
$1\%$ & 9  & $8.3\%$\\
$2\%$ & 9  & $15.5\%$\\
\bottomrule
\end{tabular}
\end{center}

表中没有哪个设定凭自身就是“正确答案”。它的作用是暴露结论依赖什么：参照类改变会移动先验，设备在真实部署中的性能下降会移动似然。若四种设定都支持同一行动，决定较稳健；若结论随一项合理改动而反转，就应把不确定性本身作为主要结果。贝叶斯计算使条件关系清楚，却不能把条件假设变成无条件事实\parencite{howson2006}。

\minitopic{区分风险、参数不确定与模型无知}风险是已经列出结果并能合理赋予概率，例如一枚已知公平硬币的下一次结果。参数不确定是知道模型形式，却不确定灵敏度究竟为$0.85$还是$0.90$。模型无知更深：检测可能受某种尚未列入的药物、实验室污染或人群选择机制影响。把三者都压成一个精确概率，会制造虚假精确。处理参数不确定，可以给区间、分布或多组情景；处理模型无知，则需要主动寻找遗漏假设、设置异常监测并保留暂缓判断的可能。

这一区分也限制“证据越多，置信度越高”的口号。若新增数据只是同一有偏模型的重复输出，样本量增加会让错误结论更稳定；若新证据揭示模型遗漏，合理更新可能不是在原模型中移动概率，而是重画整个假设空间。成熟的形式分析因此至少包含两轮提问：在模型内怎样更新，以及什么迹象要求更换模型。

\minitopic{从概率阈值到信息价值}设医生在患病概率超过$6.25\%$时倾向治疗，而侵入性治疗只有在概率超过$70\%$时才可接受。一次阳性后的合理范围若是$8\%$至$27\%$，它已经跨过“继续检查或低风险处置”的门槛，却没有跨过侵入性治疗门槛。此时第二项检测的价值不在于让医生“知道得更多”这一抽象事实，而在于它是否有较大机会把患者移到另一行动区间，并且这种改变的收益是否超过检测费用、等待和误差。

信息价值也可能为负。若复检与初检共享污染源，它提供的表面确认会诱发过度自信；若结果返回太晚，等待损失可能超过信息收益；若机构没有能力依据结果改变行动，检测只增加数据而不改善决定。决策分析应比较“立即行动”“先复检”“采用不同原理的检测”“暂缓并观察”等完整策略，而不是孤立比较两个概率\parencite{jeffrey1983}。

\minitopic{形式主义的最强反对与有限辩护}反对者担心，数字会把争议藏进输入：谁决定参照类、效用和允许的假设，谁就预先决定了答案。这个批评击中了把模型当作中立裁判的用法，却没有证明形式化毫无价值。恰当的回应不是声称输入完全客观，而是把输入选择、敏感性和未建模因素一并公开。与只说“综合考虑后风险较高”相比，一个可复算模型至少让批评者知道应攻击哪个基础率、哪项独立性假设或哪种价值权重。

\begin{argumentbox}
一份可问责的决策报告应分成两页。第一页是认识报告：命题、数据来源、先验、似然、依赖结构、模型替代和敏感性范围。第二页是行动报告：可选行动、结果价值、不可折算的权利约束、信息价值、复核权限和触发重新评估的条件。把两页合成一个“模型建议”，会掩盖从证据到行动之间新增的规范前提。
\end{argumentbox}

\index{模型不确定性}
\index{敏感性分析}
\index{信息价值}
